\section{Boundary rules for the first-match partition}\label{app:equality}
\begin{samepage}
All cuts in the partition apply to the original pair. The table gives
the first-match row for each equality, after all earlier rows are tested.
The notation is that of Section~\ref{sec:partition}.
\end{samepage}
\[
\begin{array}{c|c}
\text{boundary}&\text{first-match row}\\ \hline
R=Y^{5/3-\epsilon_0}&O_R\\
M=Y^{2/3-\epsilon_0}&O_M\\
\rad(M)=Y^{2/3-k_0}&O_{\rad M}\\
K_{\rm res}=Y^{\Delta_{\rm exp}}&\text{fails }Q0;\text{ continue testing}\\
L=Y^{1/3-l_0}&\text{fails }Q1;\text{ continue testing}\\
\min(J,K)=Y^{1/3-j_0}&\text{fails strict }Q2;\text{ test later rows}\\
u=t&u\text{-chart}\\
s=Y^{a_0}&3u\text{ or }3t,\text{ as selected}\\
s=Y^{b_0}&5W\\
\min(v,w)=Y^{p_0}&4E\\
\Pi=Y^{p_0},\quad v,w>Y^{p_0}&4H\\
\Pi>Y^{p_0},\quad v,w>Y^{p_0}&4P\\
\end{array}
\AutoTag{eq:old:A_1}
\]
The rows for \(4E,4H,4P\) require that all earlier predicates fail and
that \(Y^{a_0}<s<Y^{b_0}\). Equality at the lower-corner boundary fails the
strict \(Q2\) cut. It then goes to the first later row whose predicate
holds. The weak lower-corner inequalities describe the complements of the
strict \(Q2\) cuts.

The row predicates are Boolean functions of the original pair. A
counting argument may round an interval endpoint outward to an integer.
This only enlarges the counted set. No real rounding enters the
definition of a source row. Take the explicit inverse in
\eqref{eq:old:2_2}. Then test positivity, integrality, the gcd, the shell
and the mask. Each accepted reconstructed tuple then gives at most one
original pair.

Fix the ordered residual tuple that determines \(L\). Then the
identity \(ut=L\) determines the opposite shift. Hence a sum over \(u\) or
over \(t\) is a sum over one source coordinate. It is not a product
of independent shift ranges. In the \(4H\) and \(4P\) rows this identity
forces the opposite shift. In \(5W\) both shifts stay in the
original-cell source. Its geometric envelope adds no new shift sum.

The four $Q2$ count contributions fall into two populations:
\[
\begin{array}{c|c|c}
\text{population}&\text{contributions}&\text{selected chart}\\ \hline
Q2u&Q2u_{\rm int}+Q2u_{\rm harm}&u\le t\\
Q2t&Q2t_{\rm int}+Q2t_{\rm harm}&t<u
\end{array}
\]
Section~\ref{sec:ledger} computes their savings. The two terms on a
line are added to bound one chart population. They are not separate rows.
Their exponents and savings are
\[
\begin{array}{c|c|c|c}
\text{term}&\text{chart contribution}&\text{exponent}&\text{saving}\\ \hline
Q2t_{\rm int}&\text{interval term in }Q2t&1/2+6S_0&1/6-6S_0=972179/7500000\\
Q2t_{\rm harm}&\text{harmonic term in }Q2t&2/3-S_0&S_0\\
Q2u_{\rm int}&\text{interval term in }Q2u&1/2+3S_0&1/6-3S_0=2222179/15000000\\
Q2u_{\rm harm}&\text{harmonic term in }Q2u&2/3-4S_0&4S_0=30869/1250000
\end{array}
\AutoTag{eq:old:A_2}
\]
In each chart the larger exponent controls the sum. This is the smaller
saving. No term is an additional source population.

The partition is disjoint and exhaustive. Hence
\begin{align}
\#\Omega_Y={}&\#O_R+\#O_M+\#O_{\rad M}+\#Q0+\#Q1
 +\#Q2u+\#Q2t \notag\\
 &+\#3u+\#3t+\#4E+\#4H+\#4P+\#5W.
\AutoTag{eq:old:A_3}
\end{align}
The first-match convention is essential. A relaxed set that bounds one
row may overlap another relaxed counting set. Still, every original pair
goes to exactly one row. So the populations satisfy the exact
identity \eqref{eq:old:A_3}.
